%% =====================================================================
%%  Chapter 5: Computational Experiments
%%  Purpose, settings (everything needed to reproduce), results, discussion.
%%  The experiment and the numbers below are examples (the numbers are made up).
%%  Replace them with your own experiments.
%% =====================================================================
\chapter{Computational Experiments}
\label{chap:experiments}

\begin{writingnote}
This chapter is a writing and formatting example. The experimental settings and numerical values,
including means and standard deviations, are fictional, not measured results.
Replace them with your own experiments and findings.
\end{writingnote}

This chapter describes the computational experiments conducted to evaluate the proposed method.
\begin{writingnote}
First, state what the experiments are designed to examine (their purpose or research questions),
in line with the objective stated in the Introduction.
If you examine several things, present them as numbered questions or experiments,
and report the results and the discussion in the same order.
\end{writingnote}

\section{Experimental Settings}
\label{sec:setup}

\begin{writingnote}
Give enough information for others to reproduce the experiments:
the instances (data), the methods compared, the evaluation metrics, the parameters and the computing environment.
Explain why you chose the instances and the methods compared
(e.g., the instances are widely used in previous studies, or the method is a representative existing method).
Plan the experiments so that the comparison is fair; the points to check are listed as comments at this position in the source file.
\end{writingnote}

%% Points to check when planning the experiments
%% - Separate tuning/validation data from the final test data; do not tune using final test results.
%% - Use the same test instances for all methods and report budgets and stopping criteria.
%% - Run randomized methods with multiple seeds; report the seeds, instance count and runs per instance.
%% - Distinguish variation across instances from variation across seeds on the same instance.
%%   State which source of variation each mean and standard deviation summarizes, also in captions.
%% - Report parameter search ranges, the number of tuning trials and the selection criterion.

%% \paragraph{...} is a run-in heading that does not appear in the table of contents.
\paragraph{Instances}
We generated 10 instances of the 0-1 knapsack problem for each number of items $n \in \{100, 200, \dots, 1000\}$.
The values $v_i$ and weights $w_i$ are integers drawn uniformly from 1 to 100,
and the capacity $W$ is half of the total weight.

\paragraph{Methods compared}
We compared the proposed method with an exact method based on dynamic programming~\cite{kellerer2004knapsack}
and with the greedy algorithm described in \cref{sec:greedy}.

\paragraph{Evaluation metrics}
We evaluated the gap to the optimal value and the computation time.
The gap is $100 (z^{*} - z) / z^{*}$ [\%], where $z^{*}$ is the optimal value
and $z$ is the objective value of the obtained solution.

\paragraph{Parameters and environment}
%% Put tables in the table environment with the caption ABOVE the table.
%% Use \toprule, \midrule and \bottomrule (booktabs) and no vertical rules.
\Cref{tab:parameters} lists the parameters of the proposed method.
The penalty weight $\lambda$ was tuned with Optuna~\cite{akiba2019optuna}.
The methods were implemented in Python~\cite{python},
and all experiments were run on a computer with an XX CPU and XX~GB of memory\footnote{%
  Environment example: OS XX, Python XX, library XX (version XX).}.

\begin{table}[htbp]
  \centering
  \caption{Parameters of the proposed method (fictional settings)}
  \label{tab:parameters}
  \begin{tabular}{lc}
    \toprule
    Parameter & Value \\
    \midrule
    Penalty weight $\lambda$ & 10 \\
    Number of improvement iterations & 1000 \\
    Random seed & 0 \\
    \bottomrule
  \end{tabular}
\end{table}

\section{Results}
\label{sec:results}

\begin{writingnote}
Results do not speak for themselves~\cite{glasmandeal2010science}: refer to each figure and table in the text,
and state in words the overall trend and the key results that the reader should notice.
You do not need to repeat every number in a table.
Report unexpected results and results in which the proposed method is worse, too.
Before stating that one method is better than another, check that the difference is large compared with the variation
(e.g., the standard deviation).
Leave the interpretation of the results (why they were obtained) to the Discussion section.
Refer to figures and tables in the present tense (``\Cref{fig:results} shows \ldots''),
and describe what you did in the experiments in the past tense.
\end{writingnote}

%% Put figures in the figure environment with the caption BELOW the figure.
%% Store the files in figures/. Save plots as PDF (vector graphics) so that they stay sharp.
%% Place several figures side by side with the subfigure environment;
%% \cref then gives "Figure 5.1(a)".
\Cref{fig:results} shows the gap and the computation time of each method.
\Cref{fig:results-gap} shows that the gap of the proposed method is smaller than that of the greedy algorithm for every $n$.
\Cref{fig:results-time} shows that the computation time of the exact method grows rapidly with $n$,
whereas that of the proposed method grows only about linearly.

\begin{figure}[htbp]
  \centering
  \begin{subfigure}{0.48\linewidth}
    \centering
    \includegraphics[width=\linewidth]{results-gap.pdf}
    \caption{Gap to the optimal value}
    \label{fig:results-gap}
  \end{subfigure}
  \hfill
  \begin{subfigure}{0.48\linewidth}
    \centering
    \includegraphics[width=\linewidth]{results-time.pdf}
    \caption{Computation time}
    \label{fig:results-time}
  \end{subfigure}
  \caption[Performance versus the number of items (fictional data)]{Performance versus the number of items $n$ (fictional data; illustrative averages over 10 instances per $n$)}
  \label{fig:results}
\end{figure}

\Cref{tab:results} shows the results for $n = 1000$. The best value in each column is shown in bold.

\begin{table}[htbp]
  \centering
  \caption[Results for $n = 1000$ (fictional data)]{Results for $n = 1000$ (fictional means $\pm$ standard deviations for 10 instances)}
  \label{tab:results}
  \begin{tabular}{lrr}
    \toprule
    Method & Gap [\%] & Computation time [s] \\
    \midrule
    Exact (DP) & $\mathbf{0.00 \pm 0.00}$ & $83.0 \pm 21.7$ \\
    Greedy     & $1.07 \pm 0.35$ & $\mathbf{0.01 \pm 0.00}$ \\
    Proposed   & $0.22 \pm 0.08$ & $2.00 \pm 0.12$ \\
    \bottomrule
  \end{tabular}
\end{table}

\section{Discussion}
\label{sec:discussion}

\begin{writingnote}
Interpret the results; do not simply repeat them.
For each key result, explain why it was obtained.
If an explanation is a hypothesis that the experiments did not test, signal it with cautious language
(``A possible explanation for this is that \ldots'')~\cite{morleyphrasebank}.
Where possible, compare the results with those reported in previous studies, and state where they agree and where they differ.
Then discuss, with the results as evidence, under which conditions the proposed method is effective
and when it is not (its limitations).
The limitations discussed here lead to the future work in the Conclusion chapter.
\end{writingnote}

Under the conditions tested, the results indicate \ldots.
However, \ldots, which limits the applicability of the method to \ldots.
